% figures/ch02-three-point-types.tex
% 基本域中三类点的示意图：普通点、椭圆点、尖点
\begin{figure}[ht]
\centering
\begin{tikzpicture}[>=Stealth, font=\small, scale=1.3]

% --- 基本域轮廓 ---
% 两侧竖线 Re = ±1/2，底部圆弧 |τ|=1
\fill[blue!6]
  (-0.5, 3.5) -- (-0.5, 0.866) arc[start angle=120, end angle=60, radius=1]
  -- (0.5, 3.5) -- cycle;

\draw[thick, blue!60!black]
  (-0.5, 0.866) -- (-0.5, 3.5);
\draw[thick, blue!60!black]
  (0.5, 0.866) -- (0.5, 3.5);
\draw[thick, blue!60!black]
  (-0.5, 0.866) arc[start angle=120, end angle=60, radius=1];

% --- 顶部箭头：尖点∞ ---
\draw[->, thick, green!50!black, line width=1.2pt]
  (0, 3.2) -- (0, 3.8);
\node[green!50!black, right] at (0.05, 3.6) {$\infty$（尖点）};

% --- 普通点（内部随便一点） ---
\fill[gray!70] (0.15, 2.0) circle (2pt);
\node[right] at (0.18, 2.0) {$\tau_0$（普通点，局部坐标 $w = \tau - \tau_0$）};

% --- 椭圆点 i ---
\fill[red!80!black] (0, 1) circle (3pt);
\node[left, red!80!black] at (-0.05, 1) {$i$};
\node[right, font=\footnotesize] at (0.05, 0.85)
  {阶 $2$ 椭圆点，局部坐标 $w = \!\left(\dfrac{\tau-i}{\tau+i}\right)^{\!2}$};

% --- 椭圆点 ρ = e^{2πi/3} ---
\fill[blue!80!black] (-0.5, 0.866) circle (3pt);
\node[left, blue!80!black] at (-0.55, 0.866) {$\rho$};
\node[right, font=\footnotesize, blue!70!black] at (-0.42, 0.72)
  {阶 $3$ 椭圆点，局部坐标 $w = \!\left(\dfrac{\tau-\rho}{\tau-\bar\rho}\right)^{\!3}$};

% --- 标注基本域名称 ---
\node[blue!50!black] at (0, 2.6) {$\mathcal{F}$};

% --- 下方坐标轴 ---
\draw[->] (-1, 0) -- (1, 0) node[right] {$\Re(\tau)$};
\draw[->] (0, 0) -- (0, 0.3);
\node[below] at (-0.5, 0) {$-\tfrac{1}{2}$};
\node[below] at (0.5, 0) {$\tfrac{1}{2}$};
\node[below] at (0, 0) {$0$};

\end{tikzpicture}
\caption{基本域 $\mathcal{F}$ 中的三类点及其局部坐标。
普通点（灰色）的局部坐标直接用 $\tau - \tau_0$；
椭圆点 $i$（红色）有阶 $2$ 的稳定子，局部坐标需要平方；
椭圆点 $\rho$（蓝色）有阶 $3$ 的稳定子，局部坐标需要立方；
尖点 $\infty$（绿色）需要用 $q = e^{2\pi i\tau}$ 作局部坐标（在 §\ref{sec:compactification-X} 中讨论）。}
\label{fig:three-point-types}
\end{figure}
